% Copyright (C) 2017 - 2021 - Michael Baudin - Fabien Souillé

\documentclass[10pt]{beamer}

% Copyright (C) 2012 - EDF R&D - Michael Baudin

% To highlight source code
\usepackage{listings}
\definecolor{darkgreen}{rgb}{0,0.5,0}
\definecolor{violet}{rgb}{0.5,0,1}

\usepackage{algorithm,algorithmic}
\usepackage{bbm}

\usepackage{lmodern}% http://ctan.org/pkg/lm

\usetheme{Montpellier}
\setbeamertemplate{navigation symbols}{} % Remove navigation
\useoutertheme{infolines}

\usepackage[utf8]{inputenc}
\usepackage[T1]{fontenc}

%\usepackage[french]{babel}
%\uselanguage{French}
%\languagepath{French}

\def\bx{{\bf x}}
\def\RR{\mathbb{R}}

\newcommand{\pyvar}[1]{\texttt{#1}}

\def \ot {OpenTURNS}

\hypersetup{colorlinks=true}

\lstset{
  % general command to set parameter(s)
   %basicstyle=\footnotesize\ttfamily, %
   %basicstyle=\normalsize \ttfamily, %
   basicstyle=\scriptsize\ttfamily, %
   keywordstyle=\color{violet}\bfseries,%
   commentstyle=\color{darkgreen}\bfseries,%
   showspaces=false,%
   stringstyle=\color{red}\bfseries, 
   otherkeywords={NumericalMathFunction, Gumbel, Normal, TruncatedDistribution, %
       CompositeRandomVector, Uniform, PythonFunction, ComposedDistribution, RandomVector, %
	   Event, MonteCarlo, NumericalSample, IntervalMesher, Field, HistogramFactory, %
	   Graph, BoxCoxFactory, CompositeProcess, Indices, WhittleFactory, ARMAFactory, %
	   Basis, TrendFactory, BoxCoxTransform, SpatialFunction, TimeSeries, %
	   WelchFactory, GreaterOrEqual, SymbolicFunction, IndependentCopula, %
	   FORM, Cobyla, getResult, getEventProbability, run, drawImportanceFactors},
}

\newcommand{\fcar}[2] {{\mathbbm{1}}_{#1}\left(#2\right)}
\newcommand{\vect}[1]{{\underline{#1}}}
\newcommand{\idx}{\vect{\alpha}}
\newcommand{\argmin}{\operatornamewithlimits{argmin}}


\title[Water temperature metamodel]{Metamodeling the water temperature}

\author[Baudin and Souillé]{
Michaël Baudin \inst{1} \and
Fabien Souillé \inst{1} \\
\emph{with contributions by Chu Mai}
}

\institute[EDF]{
\inst{1} EDF R\&D. 6, quai Watier, 78401, Chatou Cedex - France, michael.baudin@edf.fr \and %
}

\date[]{June 2021}

%%%%%%%%%%%%%%%%%%%%%%%%%%%%%%%%%%%%%%%%%%%%%%%%%%%%%%%%%%%%%%%%%%%%%%%%%%%%%

\begin{document}

%%%%%%%%%%%%%%%%%%%%%%%%%%%%%%%%%%%%%%%%%%%%%%%%%%%%%%%%%%%%%%%%%%%%%%%%%%%%%

\begin{frame}
\titlepage
  
\begin{center}
\includegraphics[height=0.15\textheight]{figures/edf.jpg}
\end{center}

\end{frame}

\note{
Remercier Chu Mai. 
}

%%%%%%%%%%%%%%%%%%%%%%%%%%%%%%%%%%%%%%%%%%%%%%%%%%%%%%%%%%%%%%%%%%%%%%%%%%%%%

\begin{frame}
\frametitle{Contents}
\tableofcontents
\end{frame}

%%%%%%%%%%%%%%%%%%%%%%%%%%%%%%%%%%%%%%%%%%%%%%%%%%%%%%%%%%%%%%%%%%%%%%%%%%%%%
\section{Introduction}

\begin{frame}
\frametitle{Goals}

Industrial topic:
\begin{itemize}
\item Frazil ice are small crystal disks. 
\item They may appear when the water temperature falls below the fusion point. 
\item The growth can be fast. 
\item They may accumulate on water intake racks of nuclear power plants 
\cite{souille2020multi}: we focus on the Blayais power plant. 
\end{itemize}

\begin{columns}
\column{0.45\textwidth}
\includegraphics[height=0.45\textwidth]{figures/Frazil_ice_in_Yosemite_Creek.png} \\
\emph{Frazil ice in Yosemite Creek, \cite{wikiFrazil}}

\column{0.45\textwidth}
\begin{center}
\includegraphics[height=0.45\textwidth]{figures/Blayais-CNPE.PNG} \\
\emph{Blayais CNPE, © OpenStreetMap contributors}
\end{center}
\end{columns}

\end{frame}

%%%%%%%%%%%%%%%%%%%%%%%%%%%%%%%%%%%%%%%%%%%%%%%%%%%%%%%%%%%%%%%%%%%%%%%%%%%%%

\begin{frame}
\frametitle{Goals}

Goals of the study :
\begin{itemize}
\item We have measures of water temperature, air temperature, 
water level and flowrate of the Gironde river in the estuary.

\item We want to predict the frazil near the Blayais nuclear power plant. 

\item We make the hypothesis that the frazil appears when the water temperature 
of the Gironde falls below 0$\degree$ C.

\item The simulation horizon is 7 days:
we make the hypothesis that the water temperature is known from days 
1 to 6 and want to predict the water temperature on day 7 
(we may want to update this time interval). 
\end{itemize}

\end{frame}

%%%%%%%%%%%%%%%%%%%%%%%%%%%%%%%%%%%%%%%%%%%%%%%%%%%%%%%%%%%%%%%%%%%%%%%%%%%%%

\section{Data}
\begin{frame}
\frametitle{Data}

\begin{itemize}
\item We have water observations from the Magest station from 2005 to 2016 
\cite{Schmidt2014}.
\item We also have observations from the power plant water temperatures sensors 
from the Blayais power plant for 2015, 2016 and 2017.
\item Some observations are lacking for short or longer times (from minutes to 
months). 
\end{itemize}

\end{frame}

%%%%%%%%%%%%%%%%%%%%%%%%%%%%%%%%%%%%%%%%%%%%%%%%%%%%%%%%%%%%%%%%%%%%%%%%%%%%%

\begin{frame}
\frametitle{Data}

\begin{center}
\includegraphics[width=0.7\textwidth]{figures/Teau-semaines-continues.png}
\end{center}

Observations of the water temperature from the Magest station and Blayais 
power plan, merged from 2005 to 2019 at hour time step : 468 continuous 
weeks. 

\end{frame}

%%%%%%%%%%%%%%%%%%%%%%%%%%%%%%%%%%%%%%%%%%%%%%%%%%%%%%%%%%%%%%%%%%%%%%%%%%%%%

\begin{frame}
\frametitle{Data}

\begin{center}
\includegraphics[width=0.7\textwidth]{figures/Teau-semaines.png}
\end{center}

Observations of the water temperature from the Magest station and Blayais 
power plan, merged from 2005 to 2019 at hour time step : 468 continuous 
weeks. 
The vertical black line represents the cut between the days 1 to 6 (input) and 
the 7th day (output to be predicted).

\end{frame}

%%%%%%%%%%%%%%%%%%%%%%%%%%%%%%%%%%%%%%%%%%%%%%%%%%%%%%%%%%%%%%%%%%%%%%%%%%%%%

\section{Mathematical problem}

\begin{frame}
\frametitle{Mathematical problem}

Below are the goals and constraints:
\begin{enumerate}
\item Predict $T_{\textrm{future water}}$ depending on  $T_{\textrm{past water}}$. 

\item Using directly the temporal grid in input and output leads to a too 
large number of input ($24 \times 6 = 144$ time steps) and output variables (24 time 
steps). 

\item We would like to change easily the time grid. 

\item The link between input and output is complicated and not necessarily 
linear.

\item Possible to take into account other physical variables such as 
air temperature or water level (not necessarily discretised on the same grid).

\item The observations are noisy and the sample set is large (hundreds). 

\item We would like to estimate a confidence interval of the future temperature.
\end{enumerate}
\end{frame}

%%%%%%%%%%%%%%%%%%%%%%%%%%%%%%%%%%%%%%%%%%%%%%%%%%%%%%%%%%%%%%%%%%%%%%%%%%%%%

\begin{frame}
\frametitle{Suggested solutions}

Suggested solutions :
\begin{enumerate}
\item Reduce the dimension using the Karhunen-Loève decomposition (e.g. 
\cite{Sullivan2015}, chap.11).

\item Use sparse generalized polynomial chaos metamodel 
\cite{ghanem2003stochastic, xiu2002wiener, LeMaitre2010, blatman2011adaptive}.

\item The method suggested in \cite{Dutfoy2017} combine both previous methods 
(and discusses sensitivity analysis). 
\end{enumerate}

\end{frame}

%%%%%%%%%%%%%%%%%%%%%%%%%%%%%%%%%%%%%%%%%%%%%%%%%%%%%%%%%%%%%%%%%%%%%%%%%%%%%

\section{Karhunen-Loève decomposition}

\begin{frame}
\frametitle{Karhunen-Loève decomposition}
The Karhunen-Loève decomposition of a stochastic process $X(t)$ is:
$$
X(t) = \sum_{k\geq 0} \sqrt{\lambda_k} \xi_k \bphi_k(t)
$$
where $\xi_k \in\RR$ is the k-th coefficient in the decomposition:
$$
\xi_k = \frac{1}{\sqrt{\lambda_k}} \int_\rD X(t) \bphi_k(t) dt
$$
for $k\geq 0$ and
\begin{itemize}
\item $(\lambda_k)_{k\geq 0}$ are the eigenvalues, 
\item $(\bphi_k)_{k\geq 0}$ are the eigenvectors,
\end{itemize}
of the Fredholm problem. 
\end{frame}

%%%%%%%%%%%%%%%%%%%%%%%%%%%%%%%%%%%%%%%%%%%%%%%%%%%%%%%%%%%%%%%%%%%%%%%%%%%%%

\begin{frame}
\frametitle{Karhunen-Loève decomposition}

\begin{center}
\includegraphics[height=0.7\textheight]{figures/decompositionKL-1H-log.pdf}
\end{center}

Decreasing sequence of eigenvalues of the Karhunen-Loève decomposition 
of the Gironde water temperature time series. 

\end{frame}

%%%%%%%%%%%%%%%%%%%%%%%%%%%%%%%%%%%%%%%%%%%%%%%%%%%%%%%%%%%%%%%%%%%%%%%%%%%%%

\begin{frame}
\frametitle{Karhunen-Loève decomposition}

\begin{center}
\includegraphics[height=0.6\textheight]{figures/decompositionKL-1H-comparaison-249.pdf}
\end{center}

Comparison of the 249th trajectory of the Gironde water temperature 
and its projection on the truncated Karhunen-Loève decomposition with 9 modes. 

\end{frame}

%%%%%%%%%%%%%%%%%%%%%%%%%%%%%%%%%%%%%%%%%%%%%%%%%%%%%%%%%%%%%%%%%%%%%%%%%%%%%

\section{Field metamodel}

%\begin{frame}
%\frametitle{Field metamodel}

%\begin{itemize}
%\item On output, we predict the water temperature of the 7th day. 
%\item On input, we consider the water temperature for days 1 to 6. 
%\item Optionally, we want to consider the air temperature from weather predictions. 
%\item Since these are stochastic processes, using the time grid directly 
%would lead to a too large number of input and output variables: we 
%reduced the dimension of the problem with Karhunen-Loève decomposition. 
%\end{itemize}

%\end{frame}

%%%%%%%%%%%%%%%%%%%%%%%%%%%%%%%%%%%%%%%%%%%%%%%%%%%%%%%%%%%%%%%%%%%%%%%%%%%%%

\begin{frame}
\frametitle{Field metamodel}

\begin{center}
\includegraphics[height=0.3\textheight]{figures/metamodele-champ.pdf}
\end{center}

Dimension reduction and metamodel of the stochastic process for water temperature 
of the Gironde. 

\begin{itemize}
\item The function $g$ is the model from the input (past) water temperature $X$ 
to the output (future) water temperature $Y$. 
\item The projection of $X$ on the K.L. basis leads to the coefficients 
$(\xi_k)_{k\geq 0}$ and the projection of $Y$ leads to the coefficients 
$(\xi_k')_{k\geq 0}$. 
\item The model $g_\xi$ maps the K.L. coefficients.
\item The metamodel $\tilde{g}_\xi$ is a surrogate for $g_\xi$.
\end{itemize}

\end{frame}

%%%%%%%%%%%%%%%%%%%%%%%%%%%%%%%%%%%%%%%%%%%%%%%%%%%%%%%%%%%%%%%%%%%%%%%%%%%%%

\begin{frame}
\frametitle{Field metamodel: polynomial chaos decomposition}

The method has three steps: project, predict, lift.
\begin{enumerate}
\item Project. We decompose the input on $m_X \in \NN$ K.-L. modes 
$$
\xi_k(\omega) = \frac{1}{\sqrt{\lambda_k}} \int_{\mathcal{D}_{past}} X(t, \omega) \bphi_k(t) dt
$$
for $k = 1, \ldots, m_X$, where $\mathcal{D}_{past} = [0, t_{present}]$.

\item Predict. Map from the input to the output coefficients of the basis, based 
on a vector-to-vector metamodel :
$$
\boldsymbol{\xi}'(\omega) = \tilde{g}_{\boldsymbol{\xi}} (\boldsymbol{\xi}(\omega))
$$
where $\boldsymbol{\xi} \in \RR^{m_X}$, $\boldsymbol{\xi}' \in \RR^{m_Y}$. 

\item Lift. Compute the output on $m_Y \in \NN$ K.-L. modes :
$$
Y(t, \omega) = \sum_{k = 0}^{m_Y} \sqrt{\lambda^{(Y)}_k} \xi_k'(\omega) \bphi^{(Y)}_k(t)
$$
for any $t \in \mathcal{D}_{future} = [t_{present}, t_{final}]$.

\end{enumerate}

\end{frame}

%%%%%%%%%%%%%%%%%%%%%%%%%%%%%%%%%%%%%%%%%%%%%%%%%%%%%%%%%%%%%%%%%%%%%%%%%%%%%

\begin{frame}
\frametitle{Field metamodel: KL decomposition of inputs and ouputs}

We consider a basis of orthogonal multivariate polynomials $\psi_j$, 
that is, for which the weighted $\mathcal{L}^2$ scalar product is 
zero for different polynomials:
$$
\left( \psi_j,\psi_k \right) = 0,
$$
for $j\neq k$.

The output coefficients in the K.L. decomposition 
can be decomposed into a linear combination of multivariate polynomials:
$$
\boldsymbol{\xi}' 
= \sum_{j = 0}^p \boldsymbol{a}_j\psi_j \left(\mathring{\boldsymbol{\xi}} \right),
\qquad \mathring{\boldsymbol{\xi}} = T^{-1}(\boldsymbol{\xi})
$$
where $p \in \mathbb{N}$ is the number of coefficients in the truncated polynomial 
chaos decomposition, $\boldsymbol{a}_1, \ldots, \boldsymbol{a}_p \in\mathbb{R}^{m_Y}$ are the coefficients of the 
decomposition, $m_X \in \NN$ is the number of input K.-L. modes,  
$\mathring{\boldsymbol{\xi}} \in\mathbb{R}^{m_X}$ 
is a random vector, $m_Y \in \NN$ is the number of output K.-L. modes,  
$\boldsymbol{\xi}' \in \RR^{m_Y}$ 
is the output random vector, 
$T$ is the isoprobabilistic transformation, 
$\boldsymbol{\xi} \in \RR^{m_X}$ are the input K.L. coefficients. 
\end{frame}

%%%%%%%%%%%%%%%%%%%%%%%%%%%%%%%%%%%%%%%%%%%%%%%%%%%%%%%%%%%%%%%%%%%%%%%%%%%%%

\begin{frame}
\frametitle{Field metamodel}

\begin{itemize}
\item We use $m_X = 8$ K.L. coefficients on input and $m_Y = 4$ K.L. coefficients on output.

\item The distribution of the input is modeled by histograms, with independent 
marginals. 

\item The multivariate polynomial basis was based on linear enumeration. 

\item The polynomial was created based on a degree 3 sparse least squares 
polynomial : 31 coefficients to estimate. 

\item The cross validation (75\% for training) leads to a predictivity 
coefficient equal to 99\% for the first output K.L. mode and 
between 49\% and 62\% for modes 2, 3, 4. 

\item The standard deviation of the metamodel residual is an increasing function 
of time (remote future predictions are less accurate) and between 0.17 
and 0.38~$\degree$C. 

\item Some trajectories, especially with late change of trend, are 
predicted less accurately. 
\end{itemize}
\end{frame}

%%%%%%%%%%%%%%%%%%%%%%%%%%%%%%%%%%%%%%%%%%%%%%%%%%%%%%%%%%%%%%%%%%%%%%%%%%%%%

\begin{frame}
\frametitle{Field metamodel}

\begin{center}
\includegraphics[width=0.7\textwidth]{figures/validation-chaos-observation-vs-prediction.pdf}
\end{center}

For a given trajectory in the validation sub-sample, comparison with the 
trajectory predicted by the metamodel.

\end{frame}

%%%%%%%%%%%%%%%%%%%%%%%%%%%%%%%%%%%%%%%%%%%%%%%%%%%%%%%%%%%%%%%%%%%%%%%%%%%%%

\begin{frame}
\frametitle{Conclusions and future works}

OpenTURNS classes used:
\begin{itemize}
\item This work was performed using OpenTURNS v1.16. 
\item We used the \pyvar{KarhunenLoeveSVDAlgorithm} for the K.L. decomposition 
and \pyvar{FunctionalChaosAlgorithm} for P.C. decomposition. 
\item The example \cite{MetamodelFieldOT} showing how to create a metamodel 
of a field function was used as a starting point. 
\end{itemize}

OpenTURNS evolutions: development of new classes:
\begin{itemize}
\item new field classes: 
\pyvar{VertexValuePointToFieldFunction}, \pyvar{KarhunenLoeveReduction}, 
\pyvar{KarhunenLoeveValidation}. 
\item new methods for K.L. validation: \pyvar{KarhunenLoeveResult.drawEigenValues()} 
and \pyvar{ProcessSample.drawCorrelation()}
\item new metamodel classes: \pyvar{KFoldSplitter}, \pyvar{LeaveOneOutSplitter}
\end{itemize}


\end{frame}

%%%%%%%%%%%%%%%%%%%%%%%%%%%%%%%%%%%%%%%%%%%%%%%%%%%%%%%%%%%%%%%%%%%%%%%%%%%%%

\begin{frame}
\frametitle{End}

Future work:
\begin{itemize}
\item Provide the \pyvar{FieldMetamodelAlgorithm}, 
\pyvar{FieldMetamodelResult} and \pyvar{FieldMetamodelValidation} 
classes in OpenTURNS?
\item Analysis of a simplified ordinary differential equation which models 
dynamics of frazil ice concentration. 
\item Estimate confidence intervals for the temperature mean and 
of a future observation. 
\end{itemize}

\hspace{1cm}

\begin{center}
Thank you for your attention!

Questions?
\end{center}

\end{frame}

%%%%%%%%%%%%%%%%%%%%%%%%%%%%%%%%%%%%%%%%%%%%%%%%%%%%%%%%%%%%%%%%%%%%%%%%%%%%%

\begin{frame}
\frametitle{OpenTURNS: \url{www.openturns.org}}

\begin{center}
\includegraphics[width=0.8\textwidth]{figures/OT.pdf}
\end{center}
	
\begin{itemize}
\item Multivariate probabilistic modeling including dependence
\item Numerical tools dedicated to the treatment of uncertainties
\item Generic coupling to any type of physical model
\item Open source, LGPL licensed, C++/Python library
\end{itemize}


\end{frame}

\note{
OpenTURNS is a software for uncertainty quantification, uncertainty propagation, 
sensitivity analysis and metamodeling. 

The software is available with the open source LGPL licence on Linux and Windows. 

In order to use OpenTURNS, you can use directly the C++ library, or 
program your Python scripts. 
}

%%%%%%%%%%%%%%%%%%%%%%%%%%%%%%%%%%%%
\section{Références}
\begin{frame}[allowframebreaks]
\frametitle{Références}
\nocite{*}
\bibliographystyle{apalike}
\bibliography{Biblio}
\end{frame}
%%%%%%%%%%%%%%%%%%%%%%%%%%%%%%%%%%%%

\end{document}
